\section{Additional External-Backend Diagnostics}
\label{app:external_diagnostics}

The external pollution study compares an analytic advected Gaussian probe with OCPNet-3D in Dam, OpenWater, PierHarbor, and SimpleUnderwater. Each combination of scene, model, and current setting uses three localization episodes; Table~\ref{tab:app_pollution_localization} therefore aggregates 12 episodes per row.

\begin{table}[t]
\centering
\scriptsize
\setlength{\tabcolsep}{2pt}
\renewcommand{\arraystretch}{1.10}
\begin{tabularx}{\columnwidth}{@{}l l *{4}{Y}@{}}
\toprule
\textbf{Scale} & \textbf{Probe} & \textbf{$N$} & \makecell{\textbf{SR}\\\textbf{(\%)}} & \makecell{\textbf{Error}\\\textbf{(m)}} & \makecell{\textbf{$\Delta$ error}\\\textbf{(m)}} \\
\midrule
1.0 & Analytic & 12 & 41.7 & 25.93 & 0.00 \\
1.0 & OCPNet-3D & 12 & 50.0 & 20.81 & $-5.12$ \\
2.5 & Analytic & 12 & 50.0 & 26.00 & 0.00 \\
2.5 & OCPNet-3D & 12 & 50.0 & 22.64 & $-3.36$ \\
\bottomrule
\end{tabularx}
\caption{Pollution-localization diagnostic over four scenes (12 episodes per row). $\Delta$ error is measured against the analytic probe at the same current scale.}
\label{tab:app_pollution_localization}
\end{table}

OCPNet-3D lowers mean localization error by 5.12\,m at nominal current and 3.36\,m under stronger current, while matching or improving aggregate success. Containment quantities are reported separately for each field model: the analytic field tracks a remaining-mass state, while OCPNet tracks grid mass. Table~\ref{tab:app_external_scaling} reports removed mass for the analytic-field team-size sweep.

\subsection{Scene-backed cleanup scaling}

\begin{table}[!ht]
\centering
\scriptsize
\setlength{\tabcolsep}{5pt}
\renewcommand{\arraystretch}{1.08}
\begin{tabular}{r r r r}
\toprule
\textbf{Agents} & \textbf{Episodes} & \textbf{SR (\%)} & \makecell{\textbf{Removed mass}\\\textbf{(kg)}} \\
\midrule
2 & 20 & 40.0 & 0.2081 \\
4 & 20 & 100.0 & 0.4302 \\
8 & 20 & 100.0 & 0.4817 \\
10 & 20 & 100.0 & 0.4886 \\
\bottomrule
\end{tabular}
\caption{Analytic-containment scaling at nominal current, aggregated over OpenWater and PierHarbor with ten episodes per scene and team size.}
\label{tab:app_external_scaling}
\end{table}

Increasing the team from two to four agents raises aggregate success from 40\% to 100\% and approximately doubles removed mass. The smaller gains from four to ten agents indicate saturation under the fixed 15\,s containment horizon.

The route and pollution results show variation across underwater scenes. Each external-backend run records the scene, current scale, dataset clock, controller settings, and observation model.
